@@M0@@, we have @@M1@@), then the same is true of @@M2@@.

This last remark applies when @@M3@@ is the quotient of @@M4@@ by a regular ideal, since we can then take for @@M5@@ the Koszul complex associated with a regular homomorphism defining the ideal of @@M6@@.

\subsection*{2.3.}
The @@M7@@-algebra structure defined in 2.1 on @@M8@@ is functorial in @@M9@@ and @@M10@@. We can therefore globalize the preceding constructions: let @@M11@@ be a scheme, @@M12@@, @@M13@@ two @@M14@@-schemes; we can then define the morphisms (2.1.1) locally on @@M15@@; they glue by functoriality and define on @@M16@@ an associative graded @@M17@@-Algebra structure.

\sgaStableTarget{sga6.e07.lem.2-4}\textbf{Lemma 2.4.} \uline{Let} @@M18@@ \uline{be two} @@M19@@-\uline{schemes}, \uline{and} @@M20@@ \uline{a flat} @@M21@@-\uline{Module}. \uline{Then for every} @@M22@@, \uline{we have a canonical isomorphism}:
\sgaStableTarget{sga6.e07.eq.2-4-1}@@M23@@

It suffices to define a functorial isomorphism in the affine case; we can then define it globally by gluing in the general case.

Let @@M24@@, @@M25@@, @@M26@@; @@M27@@ is defined by a flat @@M28@@-module @@M29@@. Let @@M30@@ be a projective resolution of @@M31@@ over @@M32@@; the @@M33@@ are the homology groups of @@M34@@. Since @@M35@@ is flat over @@M36@@, we have a canonical isomorphism: @@M37@@. Furthermore, the canonical isomorphism: @@M38@@ gives a canonical isomorphism: @@M39@@, whence the result.

\sgaStableTarget{sga6.e07.prop.2-5}\textbf{Proposition 2.5.} \uline{Let} @@M40@@ \uline{be a scheme}, @@M41@@ \uline{and} @@M42@@ \uline{two closed subschemes of} @@M43@@, @@M44@@. \uline{Suppose that for every} @@M45@@, \uline{there exists an open neighborhood}
